% ==============================================================================
% Fichier : front/03_ndap_sukulu.tex
% Titre   : NDAP SUKULU, l'environnement compagnon du cours
% Auteur  : MINKA MI NGUIDJOÏ Thierry Emmanuel
% Version : 1.0 -- novembre 2026
% ==============================================================================
% Texte commun aux cours de l'auteur ; seul le tableau des labs est propre au
% cours (fichier front/03b_labs.tex).
% ==============================================================================

\chapter*{NDAP SUKULU, la salle de classe compagnon}
\addcontentsline{toc}{chapter}{NDAP SUKULU, la salle de classe compagnon}
\markboth{NDAP SUKULU}{NDAP SUKULU}

\begin{center}
    \includegraphics[width=0.62\textwidth]{figures/logos/logo_ndap_sukulu.png}
\end{center}

\section*{Un nom venu du pays bassa}

\begin{quote}\itshape
Le soir tombe sur Éséka. Un enfant revient des champs avec son grand-père et
lui demande pourquoi il doit marcher si loin, chaque matin, pour aller à
l'école. Le vieil homme s'arrête au bord du chemin et répond, en bassa~:
« \textup{Ndap sukulu}, la salle de classe, ce ne sont pas les murs. Ce sont
les murs qui viennent après. La salle de classe, c'est l'endroit où quelqu'un
t'attend pour t'apprendre ce qu'il sait. » L'enfant insiste~: « Et qui
m'attend~? » Le grand-père sourit~: « \textup{Malet yon}. Ton enseignant. »
\end{quote}

En bassa, la langue paternelle de l'auteur, \textbf{\textit{ndap sukulu}}
désigne la salle de classe et \textbf{\textit{malet yon}} signifie « ton
enseignant ». L'environnement compagnon de ce cours porte ce nom parce qu'il
prolonge la salle de classe au-delà des heures de cours~: il attend
l'étudiant le soir, le week-end, à distance, avec des questions, des pièces à
examiner et une correction immédiate. Il est édité par \textbf{MaletYon},
« ton enseignant », et mis à la disposition des étudiants et des praticiens
qui souhaitent se perfectionner (\href{mailto:maletyon@proton.me}{maletyon@proton.me}).

\section*{Ce que l'on y trouve}

\begin{itemize}
    \item \textbf{Une fiche d'auto-évaluation par chapitre.} Chaque fiche
          puise vingt questions dans une banque de quarante (QCM, vrai ou
          faux, textes à trous). Elle dure trente minutes~; la copie validée
          est verrouillée et notée aussitôt. Le seuil de réussite est de
          60~\%.
    \item \textbf{Des laboratoires pratiques.} Chaque lab place l'étudiant
          devant les pièces d'une affaire du fil rouge (tableurs, journaux,
          contrats, procès-verbaux) dans un bureau d'analyste virtuel,
          ouvert dans le navigateur. Il dure soixante minutes, se termine
          par des points de contrôle notés automatiquement et renvoie aux
          sections du livre qu'il mobilise.
    \item \textbf{Un suivi pour l'enseignant}~: rapport de classe, notes par
          fiche et par lab, journal d'activité.
\end{itemize}

Les fiches et les labs sont des travaux à faire hors des heures de cours.
Ils ne remplacent ni le livre ni les exercices de fin de chapitre~: ils en
vérifient la maîtrise.

\section*{Comment y travailler}

\begin{enumerate}
    \item Se connecter au portail NDAP SUKULU avec ses identifiants et le
          code de cours \textbf{\codeUE}.
    \item Pour un lab, commencer par vérifier l'empreinte SHA-256 des pièces
          (\texttt{sha256sum -c MANIFESTE\_SHA256.txt}), puis travailler sur
          la copie du dossier \texttt{travail}~: les pièces originales sont
          en lecture seule et l'activité est journalisée, comme dans une
          vraie mission.
    \item Lire le rappel de cours affiché par le lab et relire les sections
          citées avant de répondre.
\end{enumerate}

\begin{boxinfo}
\textbf{Si le portail n'est pas encore ouvert.} En attendant le déploiement
de NDAP SUKULU, les fiches de questions sont fournies au format texte dans le
dossier \texttt{fiches} du cours, prêtes à être distribuées. Les corrections
des fiches figurent dans le livret de correction, à la suite des corrigés des
exercices de chaque chapitre.
\end{boxinfo}

\section*{Les laboratoires du cours}

% Tableau des labs du cours SIF (SEC3266) dans NDAP SUKULU (propre au cours).
Le cours compte dix laboratoires, disponibles sous le code \codeUE. Le
dernier reconstitue l'attaque de l'Opération GAROUA à partir de ses
journaux.

\begin{tableaucours}{L{2.4cm}L{6.4cm}L{1.6cm}L{1.8cm}}{%
    \tableauHeader{Lab} &
    \tableauHeader{Sujet} &
    \tableauHeader{Chapitre} &
    \tableauHeader{Opération}
}
SIF-LAB-01 & Empreintes et HMAC~: retrouver les fichiers modifiés &
\ref{chap:crypto-sym} & GAROUA \\
SIF-LAB-02 & Certificats~: lire une chaîne, repérer les défauts &
\ref{chap:crypto-asym} & LIZBIZ \\
SIF-LAB-03 & Analyse de risques EBIOS simplifiée & \ref{chap:risques} &
LIZBIZ \\
SIF-LAB-04 & Classer les actifs et coter les besoins & \ref{chap:actifs} &
LIZBIZ \\
SIF-LAB-05 & Capture réseau~: clair, balayage, ARP & \ref{chap:reseaux1} &
GAROUA \\
SIF-LAB-06 & Règles de pare-feu & \ref{chap:reseaux2} & GAROUA \\
SIF-LAB-07 & Revue des comptes & \ref{chap:iam} & GAROUA \\
SIF-LAB-08 & Durcissement Linux & \ref{chap:systemes} & LIZBIZ \\
SIF-LAB-09 & Revue de code & \ref{chap:applications} & MEFOU \\
SIF-LAB-10 & Journaux~: reconstituer l'attaque & \ref{chap:operations} &
GAROUA \\
\end{tableaucours}
\vspace{-0.8em}
\begin{center}{\small\itshape Les laboratoires du cours dans NDAP SUKULU}\end{center}

